\documentclass[11pt,a4paper]{report}

% Packages
\usepackage[utf8]{inputenc}
\usepackage[T1]{fontenc}
\usepackage{graphicx}
\usepackage{hyperref}
\usepackage{booktabs}
\usepackage{tabularx}
\usepackage{listings}
\usepackage{amsmath}
\usepackage[margin=1.27cm]{geometry}
\usepackage[authoryear,round]{natbib}
\usepackage{url}
\usepackage{xcolor}
\usepackage{float}
\usepackage{titlesec}
\usepackage{setspace}

% Formatting
\singlespacing
\hypersetup{
    colorlinks=true,
    linkcolor=blue,
    filecolor=magenta,      
    urlcolor=cyan,
    citecolor=blue,
}

% Title Page Info
\title{TCM-Sage: A Glass-Box Evidence-Synthesis System for Classical Traditional Chinese Medicine via Hybrid RAG and Knowledge Graph Integration}
\author{ZHENG Zian, Andy (Student ID: 22231153) \\ 
        Department of Computer Science \\ 
        Hong Kong Baptist University \\
        Computing and Software Technologies (CST) Concentration}
\date{April 8, 2026}

\begin{document}

\maketitle

% Declaration Page
\newpage
\thispagestyle{empty}
\begin{center}
\vspace*{5cm}
{\Large\bfseries Declaration}
\end{center}
\vspace{2cm}
I hereby declare that all the work done in this Final Year Project is of my independent effort. I also certify that I have never submitted the idea and product of this Final Year Project for academic or employment credits.

\vspace{3cm}
\noindent
\begin{tabular}{ll}
Signature: & \underline{\hspace{6cm}} \\[1cm]
Name: & ZHENG Zian (Andy) \\[1cm]
Date: & April 8, 2026 \\
\end{tabular}
\newpage

\maketitle

\begin{abstract}
Traditional Chinese Medicine (TCM) practitioners face significant challenges in retrieving and synthesizing evidence from a vast corpus of classical literature. Existing Large Language Models (LLMs) often function as ``black boxes,'' providing unverifiable claims that lack the precision required for clinical or academic reference. This project introduces TCM-Sage, an evidence-synthesis tool designed to provide transparent, citation-backed insights from 17 foundational TCM texts comprising 3.72 million characters and 12,204 indexed chunks. By implementing a Retrieval-Augmented Generation (RAG) pipeline with hybrid retrieval---combining DashScope text-embedding-v4 vector search with the SymMap v2.0 knowledge graph (18,450 entities)---TCM-Sage ensures every claim is traceable to its original source. Evaluation via a blind A/B Arena system with 56 votes from TCM students and practitioners yielded a paired t-test result of $t = 3.30$, $p = 0.0018$ ($p < 0.01$), and Cohen's $d = 0.47$ (medium effect), demonstrating statistically significant preference for the RAG system over a general-purpose LLM with web search.
\end{abstract}

\tableofcontents

\chapter{Introduction}

\section{Background and Motivation}
Traditional Chinese Medicine (TCM) practitioners rely on a massive body of classical literature to guide clinical reasoning and theoretical understanding. This corpus, spanning over two millennia, includes 17 foundational texts comprising millions of characters written in classical Chinese. The complexity of this language, combined with the sheer volume of data, makes manual evidence retrieval and cross-text synthesis an arduous task for even experienced practitioners. Modern TCM education and practice require tools that can efficiently navigate these texts to provide evidence-based support.

The motivation for this project stems from both personal and professional observations. My father first suggested exploring the intersection of AI and medicine, noting the potential for technology to assist in specialized fields. This direction was further reinforced through discussions with Kenny Woo Shi Nam, a Year 5 student at the HKBU School of Chinese Medicine with a 3.97 cGPA, currently interning at the Guangdong Provincial Hospital of TCM. Kenny highlighted the critical need for a tool that could provide reliable citations from classical texts, as current digital resources are often fragmented or lack the depth required for serious study.

Originally, my Final Year Project (FYP) topic was an AI-based course selection system. However, recognizing the greater impact and technical challenge of the TCM domain, I proposed a shift to building a TCM AI tool. This change was approved by my supervisor, Dr. Zhang Ce, who recognized the value of applying AI to preserve and utilize cultural heritage through modern engineering.

The core technical premise of TCM-Sage is that AI can serve as an external memory bank. While human practitioners may struggle to memorize every clause across 17 different texts, a Retrieval-Augmented Generation (RAG) system can store and retrieve this information with perfect recall. By integrating these texts into a structured digital framework, we can provide a level of knowledge synthesis that exceeds human capacity while maintaining the rigor of the original sources.

\begin{table}[H]
\centering
\caption{The 17 Classical TCM Texts Included in TCM-Sage}
\begin{tabularx}{\textwidth}{lX}
\toprule
\textbf{Category} & \textbf{Included Texts} \\
\midrule
Foundational Theory & Huangdi Neijing (Suwen, Lingshu), Nanjing \\
Acupuncture & Zhenjiu Jiayijing \\
Shanghan/Wenbing & Shanghan Lun, Jingui Yaolue, Wenbing Tiaobian, Wenre Lun \\
Herbology/Formulary & Shennong Bencao Jing, Bencao Gangmu \\
Jin-Yuan Four Masters & Suwen Xuanji Yuanbing Shi, Rumen Shiqin, Piwei Lun, Lanshi Micang, Danxi Xinfa, Gezhi Yulun \\
General Clinical & Beiji Qianjin Yaofang, Neiwaishang Bianhuo Lun \\
\bottomrule
\end{tabularx}
\label{tab:texts}
\end{table}

\section{Problem Statement}
The primary obstacle to using AI in TCM is the "Black Box" nature of existing general-purpose Large Language Models (LLMs) like ChatGPT or Qwen. When a user queries these models for TCM knowledge, the system provides a response without any verifiable link to the source material. This lack of transparency makes it impossible for practitioners to confirm whether a claim is based on a classical text or is a fabrication of the model's internal weights.

The core issue is not merely hallucination, but unverifiability. While all AI systems, including TCM-Sage, may initially produce errors, the critical difference is that general LLMs offer no mechanism for verification. For example, a TCM student testing ChatGPT found that the model confused Clause 82 of the \textit{Shanghan Lun} with Clause 2. The model confidently presented the wrong content as if it were an exact quote, which could lead to incorrect clinical reasoning if taken at face value.

Existing TCM AI products, such as HuatuoGPT, Qihuang, and Xunfei, have largely pursued an "AI Doctor" approach. These tools focus on providing direct diagnoses and prescriptions. However, this approach has seen low clinical practicality, with usage rates for tools like Qihuang reported at less than 15\%. Several factors contribute to this failure:
\begin{enumerate}
    \item \textbf{Data Sovereignty:} Hospitals and practitioners are often unwilling to upload sensitive patient data to third-party cloud-based AI systems.
    \item \textbf{Clinical Risk:} Direct diagnostic suggestions carry high liability and often lack the nuanced reasoning required for TCM's "Bian Zheng Lun Zhi" (Pattern Differentiation and Treatment).
    \item \textbf{Lack of Evidence:} None of these tools focus on the synthesis of evidence from classical texts with traceable, clause-level citations.
\end{enumerate}

There is an apparent gap for a tool that prioritizes evidence synthesis over direct diagnosis. Practitioners need a ``Glass Box'' system where every claim is backed by a traceable citation, allowing the human expert to remain the final decision-maker while the AI handles the literature retrieval.

\section{Proposed Solution}
TCM-Sage is designed as an Evidence-Synthesis Tool rather than an AI Doctor. It adopts a "Glass Box" architecture, ensuring that every piece of information provided to the user is directly traceable to a source in the 17 classical texts. The system does not replace the practitioner; it empowers them by providing a synthesized view of the literature relevant to a specific case or query.

The core of the solution is a Retrieval-Augmented Generation (RAG) pipeline that utilizes Hybrid Retrieval. This approach combines:
\begin{itemize}
    \item \textbf{Vector Search:} Using DashScope text-embedding-v4 to find semantically similar chunks of text across the corpus.
    \item \textbf{Knowledge Graph Integration:} Utilizing the SymMap v2.0 knowledge graph to bridge the semantic gap between modern medical terms and classical Chinese descriptions.
\end{itemize}

While TCM-Sage can provide treatment suggestions when provided with full "Si Zhen" (Four Examinations) case data, these suggestions are always presented alongside the relevant literature. This design principle ensures that the system remains a reference tool, maintaining the practitioner's role in the diagnostic process while providing the necessary evidence to support their conclusions.

\begin{figure}[H]
\centering
\fbox{\parbox{0.8\textwidth}{\centering [PLACEHOLDER: System architecture overview diagram showing RAG pipeline, SymMap integration, and Citation Panel]}}
\caption{High-level architecture of the TCM-Sage system.}
\label{fig:arch}
\end{figure}

\section{Objectives}
The project aims to achieve the following technical and functional objectives:
\begin{enumerate}
    \item \textbf{Modular RAG Pipeline:} Build a specialized pipeline for classical Chinese medical texts, optimized for the unique linguistic structures of the 17 target texts.
    \item \textbf{Knowledge Graph Integration:} Integrate the SymMap v2.0 knowledge graph, which contains 18,450 entities and 21,476 relationships. This includes developing a "crosswalk bridge" to map modern symptoms to classical entities.
    \item \textbf{Clause-Level Chunking:} Implement a structured chunking strategy for key texts. For instance, the \textit{Shanghan Lun} is divided into its 388 original clauses, and the \textit{Jingui Yaolue} into 489 clauses, ensuring that citations point to specific, meaningful units of text.
    \item \textbf{Arena Evaluation System:} Develop a blind A/B evaluation system inspired by the LMSYS Chatbot Arena. This system will allow for statistical validation of the RAG output using T-Tests and Cohen's d effect size calculations.
    \item \textbf{Local Deployment Support:} Design the system architecture to support local LLM inference via Ollama or LMStudio, addressing data privacy concerns. The current deployment uses cloud-based DashScope embeddings, but the modular design allows substitution with local embedding models for fully offline operation.
    \item \textbf{Transparent Citation UI:} Provide a dedicated citation panel in the frontend that allows users to click through from a claim to the exact source text, including metadata about the text's origin and context.
\end{enumerate}

\section{Scope and Limitations}
The scope of TCM-Sage is strictly defined to ensure depth and accuracy within the project's timeframe. The system covers 17 classical TCM texts that represent the core of TCM theory and practice, including foundational theory, acupuncture, Shanghan/Wenbing, herbology, and the works of the Jin-Yuan Four Masters.

However, several areas are explicitly out of scope:
\begin{itemize}
    \item \textbf{Specialized Fields:} Surgery, orthopedics, ophthalmology, and specialized gynecology or pediatrics texts are not included.
    \item \textbf{Pulse Diagnosis:} Monographs dedicated solely to pulse diagnosis are excluded due to the difficulty of representing tactile data in a text-based RAG system.
    \item \textbf{Clinical Decision Making:} The system is not intended for direct clinical decision-making. It is an evidence reference tool.
\end{itemize}

The quality of the system's output is inherently bounded by the capabilities of the underlying LLM. While the RAG pipeline provides the correct context, the final synthesis depends on the model's ability to process classical Chinese and generate coherent English or modern Chinese responses.

\section{Report Structure}
The remainder of this report is organized as follows:
\begin{itemize}
    \item \textbf{Chapter 2: Literature Review} examines the current state of RAG technology, TCM AI systems, and the role of knowledge graphs in medical informatics.
    \item \textbf{Chapter 3: Methodology} details the design of the RAG pipeline, the integration of SymMap v2.0, and the crosswalk bridge implementation.
    \item \textbf{Chapter 4: Implementation} describes the technical stack, including the FastAPI backend, Next.js frontend, and the local deployment configuration.
    \item \textbf{Chapter 5: Evaluation} presents the results of the Arena blind A/B testing and the statistical analysis of the system's performance.
    \item \textbf{Chapter 6: Discussion} analyzes the strengths and weaknesses of the system, addressing the feedback from practitioners and the technical challenges encountered.
    \item \textbf{Chapter 7: Conclusion} summarizes the project's contributions and outlines future work for expanding the corpus and improving retrieval accuracy.
\end{itemize}

\chapter{Literature Review}

This chapter examines the theoretical foundations and existing research relevant to the development of TCM-Sage. It covers the evolution of Retrieval-Augmented Generation (RAG), the current landscape of Traditional Chinese Medicine (TCM) AI systems, the integration of Knowledge Graphs (KGs) in medical informatics, and the methodologies for evaluating large language model outputs in specialized domains.

\section{Retrieval-Augmented Generation (RAG)}

Retrieval-Augmented Generation (RAG) has emerged as a dominant architecture for addressing the inherent limitations of static Large Language Models (LLMs). Standard LLMs suffer from "knowledge cutoff" and hallucinations, where the model generates factually incorrect but plausible-sounding information. RAG mitigates these issues by combining a pre-trained generator with an external retrieval mechanism \cite{lewis2020rag}.

\subsection{RAG Fundamentals and Architecture}

The RAG process involves two primary stages: retrieval and generation. During retrieval, the system identifies relevant documents from a non-parametric knowledge base based on the user's query. These documents are then prepended to the user's prompt, providing the generator with a grounded context for synthesis. This architecture allows the model to access up-to-date or specialized information without the need for expensive retraining or fine-tuning.

\subsection{RAG vs. Fine-Tuning Trade-offs}

While fine-tuning adapts a model's internal weights to a specific domain, RAG offers several advantages for medical applications. Fine-tuning is computationally expensive and often results in "catastrophic forgetting," where the model loses general reasoning capabilities. Furthermore, fine-tuned models remain "black boxes" that cannot provide verifiable citations for their outputs.

TCM-Sage adopts a RAG-first approach for four primary reasons. First, fine-tuning on classical Chinese texts is prohibitively expensive and may yield poor results due to the scarcity of high-quality instruction-tuning pairs. Second, the RAG knowledge base can be updated or expanded by simply adding new text chunks to the vector store, requiring no retraining. Third, the modular nature of RAG allows the base LLM to be swapped as newer, more capable models become available. Finally, this architecture provides a foundation for future hybrid approaches where a fine-tuned TCM model can be paired with a RAG pipeline for maximum precision.

\subsection{Advanced RAG Techniques}

Recent advancements have moved beyond "naive" RAG toward more sophisticated pipelines. Reranking has become a critical component for improving retrieval precision. A reranker model evaluates the initial set of retrieved documents and re-orders them based on semantic relevance, ensuring that the most pertinent information appears at the top of the context window \cite{gao2024rag}.

Self-RAG and verification techniques further enhance reliability by allowing the model to critique its own retrieved context and generated output \cite{asai2024selfrag}. In domain-specific applications like medicine, these verification loops are essential for ensuring that the synthesized advice aligns with the source material. MedRAG research highlights the unique challenges of medical retrieval, including the need for high-precision terminology matching and the handling of complex clinical reasoning \cite{luo2024medrag}.

\section{TCM Knowledge Representation and Existing Systems}

Traditional Chinese Medicine presents unique challenges for digital representation and AI processing. The core of TCM knowledge resides in classical texts written over two millennia, featuring linguistic structures that differ significantly from modern Chinese.

\subsection{Classical Chinese Text Challenges}

Classical Chinese is characterized by extreme brevity and high semantic density. A single character often carries multiple meanings depending on the context (one-character-multiple-meanings). Archaic grammar and semantic drift across different dynasties further complicate automated analysis. For example, a term used in the \textit{Huangdi Neijing} (approx. 300 BCE) may have a different clinical implication in the \textit{Wenbing} literature of the Qing Dynasty (1644-1911 CE).

\subsection{Existing TCM AI Products}

Several major research institutions and corporations have developed TCM-specific AI systems. HuatuoGPT, developed by CUHK Shenzhen, achieved an 88.1\% score on the TCM licensing exam by fine-tuning on a large corpus of medical data \cite{huatuogpt2023}. Qihuang 2.0, a 320B parameter multimodal model, represents the scale-up approach to TCM AI. Corporate entries include iFlytek's TCM assistant, deployed in 107 clinics, and Alibaba Health's upcoming TCM initiatives.

\subsection{The "AI Doctor" Gap}

A critical observation of existing systems is their pursuit of the "AI Doctor" paradigm. These tools aim to provide direct diagnoses and prescriptions, effectively replacing the practitioner's reasoning. However, surveys indicate that TCM practitioners cite trust, unverifiability, and data sovereignty as primary barriers to adoption. Clinical usage of tools like Qihuang remains below 15\% because they lack the transparency required for professional medical work.

TCM-Sage identifies a significant gap: no existing tool focuses on evidence-synthesis from classical texts with traceable, clause-level citations. By prioritizing the "Glass Box" approach, TCM-Sage serves as a reference tool that supports, rather than replaces, the human expert.

\section{Knowledge Graphs in Medical AI}

Knowledge Graphs (KGs) provide a structured, symbolic representation of domain knowledge that complements the sub-symbolic nature of LLMs. In medical AI, KGs ground the model's output in established facts and relationships.

\subsection{SymMap v2.0 Integration}

SymMap v2.0 is a comprehensive TCM knowledge graph containing 18,450 entities and 21,476 relationships \cite{symmap2020}. It integrates modern pharmacology with traditional herbology, providing a bridge between different medical paradigms. TCM-Sage uses SymMap to enhance retrieval by identifying related symptoms, herbs, and formulas that may not be explicitly mentioned in the user's query but are semantically linked in the TCM theoretical framework.

\subsection{Bridging Modern and Classical Terminology}

A major challenge in TCM AI is the semantic drift between modern medical terminology and classical descriptions. A practitioner might use a modern term for a symptom that appears in the \textit{Shanghan Lun} under a different name. The integration of a KG allows for entity resolution and alias mapping, ensuring that the RAG pipeline can retrieve relevant classical evidence even when the query uses modern language.

\section{Evaluation Methodologies}

Evaluating the quality of LLM-generated medical advice is difficult due to the subjective and nuanced nature of clinical reasoning. Traditional metrics like ROUGE or BLEU are insufficient for assessing factual accuracy or clinical relevance.

\subsection{LMSYS Chatbot Arena and Blind A/B Testing}

The LMSYS Chatbot Arena has popularized blind A/B testing as a gold standard for evaluating LLM performance \cite{chatbotarena2024}. In this framework, users compare two anonymized outputs and vote for the superior response. This method captures human preferences and nuances that automated rubrics miss.

\subsection{Statistical Validation}

To ensure that the results of blind A/B testing are scientifically valid, TCM-Sage employs rigorous statistical methods. The paired T-Test is used to determine if the difference in scores between the RAG system and a baseline LLM is statistically significant. Cohen's d is calculated to measure the effect size, providing an indication of the practical significance of the improvement. These methods provide a robust mathematical foundation for claims of system superiority, moving beyond anecdotal evidence.


\chapter{System Architecture and Design}

This chapter describes the technical architecture and design principles of TCM-Sage. It details the evolution of the system through three planned phases, the construction of the specialized knowledge base, the hybrid retrieval mechanism, and the generation pipeline optimized for classical Chinese medical texts.

\section{System Overview}

TCM-Sage is a modular Retrieval-Augmented Generation (RAG) system designed for evidence synthesis from classical Traditional Chinese Medicine (TCM) literature. The high-level architecture follows a pipeline where classical texts are processed into a structured knowledge base, retrieved through a hybrid vector-graph mechanism, and synthesized by a Large Language Model (LLM) with a dedicated citation panel for verification.

\subsection{Three-Phase Evolution}

The development of TCM-Sage followed a structured, three-phase evolution planned from the project's inception. This narrative reflects a deliberate design progression rather than a reactive fix for technical failures.

\begin{itemize}
    \item \textbf{Phase 1: Naive RAG.} The initial prototype focused on a single foundational text, the \textit{Huangdi Neijing}. It utilized a lightweight all-MiniLM-L6-v2 embedding model (384 dimensions) and a basic vector search mechanism. This phase established the feasibility of retrieving classical Chinese text chunks for LLM synthesis.
    \item \textbf{Phase 2: Hybrid RAG.} Planned as the core architectural shift, Phase 2 introduced the integration of the SymMap 2.0 knowledge graph and a Self-RAG verifier. This phase moved beyond simple semantic similarity to include structured medical relationships, allowing the system to bridge modern symptoms with classical entities.
    \item \textbf{Phase 3: Production Hybrid RAG.} The final phase scaled the system to 17 classical texts and upgraded the technical stack. It implemented DashScope text-embedding-v4 (1024 dimensions) and the qwen3-rerank model for high-precision retrieval. This phase also introduced the crosswalk bridge for entity resolution and the full Arena evaluation system.
\end{itemize}

\begin{figure}[H]
    \centering
    \fbox{\parbox{0.8\textwidth}{\centering [PLACEHOLDER: System architecture diagram showing the flow from 17 texts to the final citation-backed response]}}
    \caption{High-level architecture of the TCM-Sage production system.}
    \label{fig:system_arch}
\end{figure}

\section{Knowledge Base Design}

The knowledge base is the foundation of TCM-Sage, comprising 17 classical texts with a total of 3.72 million characters. These texts are divided into 12,204 chunks using a specialized chunking strategy.

\subsection{Chunking Strategies}

TCM-Sage employs two distinct chunking strategies based on the structural characteristics of the source material. Most texts use a chapter-based approach, where chunks are created at natural paragraph or section boundaries. However, for the canonical clinical texts \textit{Shanghan Lun} and \textit{Jingui Yaolue}, a clause-level strategy is used.

\subsection{Clause-Level Chunking}

The \textit{Shanghan Lun} is divided into its 388 original clauses, and the \textit{Jingui Yaolue} into 489 clauses. Each clause represents a natural semantic unit in TCM, typically containing a diagnostic statement followed by a corresponding formula. This granularity ensures that citations point to the exact clinical rule being referenced.

To improve embedding quality, formula-aware contextual headers are prepended to each clause. For example, Clause 35 of the \textit{Shanghan Lun} is stored with the header "Shanghan Lun Di 35 Tiao Mahuang Tang:". This technique ensures that the embedding model captures both the clause number and the primary formula, significantly improving retrieval accuracy for specific clinical queries.

\begin{figure}[H]
    \centering
    \fbox{\parbox{0.8\textwidth}{\centering [PLACEHOLDER: Diagram comparing chapter-based chunking vs. clause-level chunking for Shanghan Lun]}}
    \caption{Comparison of chunking strategies for different text types.}
    \label{fig:chunking_strategy}
\end{figure}

\section{Hybrid Retrieval Architecture}

The retrieval mechanism combines dense vector search with symbolic knowledge graph traversal to achieve high recall and precision.

\subsection{Vector Search and Chronological Boost}

The system uses DashScope text-embedding-v4 (1024 dimensions) for semantic vector search. To handle the large corpus, the retriever initially fetches a broader set of candidates (k $\times$ 3) before applying a reranking step.

A unique feature of TCM-Sage is the source authority chronological boost. Distance scores are adjusted by a factor of 0.90 to 0.97 based on the historical authority of the text. This boost is topic-dependent rather than a simple linear ranking, ensuring that foundational texts like the \textit{Huangdi Neijing} are prioritized for theoretical queries, while later \textit{Wenbing} texts are favored for febrile disease queries.

\subsection{Formula-Aware Canonical Retrieval}

When the system detects a specific formula name in the user's query, it triggers a metadata-filtered search in canonical texts. This ensures that the primary source for a formula is always included in the context, even if other semantically similar but less authoritative texts are also retrieved.

\subsection{Knowledge Graph and Crosswalk Bridge}

The SymMap 2.0 knowledge graph is traversed using the NetworkX library. Entity matching is performed using the jieba segmentation tool combined with a custom alias map. The crosswalk bridge maps modern medical terms to classical entities, allowing the system to resolve queries like "hypertension" to classical patterns like "Gan Yang Shang Kang" (Liver Yang Rising).

\begin{figure}[H]
    \centering
    \fbox{\parbox{0.8\textwidth}{\centering [PLACEHOLDER: Flowchart of the hybrid retrieval pipeline showing vector search, KG traversal, and reranking]}}
    \caption{The hybrid retrieval pipeline combining vector and graph-based methods.}
    \label{fig:retrieval_flow}
\end{figure}

\section{Generation Pipeline}

The generation pipeline synthesizes the retrieved context into a coherent, evidence-backed response. It utilizes the DashScope Qwen series of LLMs, which are highly capable in processing classical Chinese.

\subsection{System Prompt and Framework}

The system prompt is structured around the "Bian Zheng Lun Zhi" (Pattern Differentiation and Treatment) framework and the "Yuan Liu" (Source and Development) methodology. It includes specific instructions for "Jun Chen Zuo Shi" (Sovereign, Minister, Assistant, and Courier) drug analysis for formulas.

\subsection{Post-Context Instruction Anchoring}

To prevent the model from ignoring formatting rules, TCM-Sage uses post-context instruction anchoring. Formatting instructions are placed after the retrieved context rather than before it. This ensures that the model's final attention is focused on the required output structure.

\subsection{Pattern Priming and User Intent}

The system employs "Pattern Priming" by providing the retrieved context in a structured markdown format (using headers and blockquotes). This encourages the LLM to produce similarly structured outputs. An ancient measurement conversion table is embedded in the prompt to ensure accurate dosage translations. Finally, a "User Intent Priority" block ensures that specific user preferences always override default formatting rules.

\section{Arena Evaluation Design}

The evaluation system is inspired by the LMSYS Chatbot Arena, providing a blind A/B testing framework for TCM practitioners.

\subsection{Blind A/B Testing}

The system compares the full TCM-Sage pipeline (RAG system) against a plain LLM baseline. The baseline model is given a generic "helpful assistant" prompt and access to a standard web search tool (DuckDuckGo). Both systems stream their responses simultaneously using Server-Sent Events (SSE), with their positions (Left vs. Right) randomly assigned to eliminate position bias.

\subsection{Statistical Validation}

The results of the arena votes are analyzed using a paired T-Test and Cohen's d effect size. These metrics provide statistical proof of the RAG system's effectiveness. The T-Test determines if the improvement is due to the system's design rather than random chance, while Cohen's d quantifies the magnitude of the improvement.

\begin{figure}[H]
    \centering
    \fbox{\parbox{0.8\textwidth}{\centering [PLACEHOLDER: Flowchart of the Arena evaluation process from query submission to statistical analysis]}}
    \caption{The blind A/B Arena evaluation and statistical validation flow.}
    \label{fig:arena_flow}
\end{figure}


\input{chapters/chapter4}

\chapter{Evaluation}
\label{ch:evaluation}

This chapter evaluates the efficacy of the TCM-Sage system compared to a non-specialized large language model (LLM) baseline. The evaluation is conducted through a blind A/B testing framework, the "Arena," involving practitioners and students from the Traditional Chinese Medicine (TCM) community.

\section{Experimental Setup}

The evaluation framework was designed to isolate the impact of the specialized RAG (Retrieval-Augmented Generation) pipeline and the TCM-specific system prompt.

\subsection{System Configurations}
Two distinct system configurations were compared:
\begin{enumerate}
    \item \textbf{Treatment (TCM-Sage):} The full pipeline described in Chapter 4, including specialized RAG context from 17 classical texts, SymMap v2.0 knowledge graph integration, and the professional TCM system prompt.
    \item \textbf{Baseline (Plain LLM):} A generic configuration using the same underlying Qwen model but with a "helpful assistant" prompt. To simulate a modern AI's capabilities, the baseline had access to \textbf{DuckDuckGo Web Search} (5 results), representing the current standard for general-purpose retrieval.
\end{enumerate}

\subsection{Blind A/B Design}
A total of 56 votes were collected from a pool of TCM students and practitioners. The evaluation used a blind design:
\begin{itemize}
    \item \textbf{Randomization:} The positions of the Treatment and Baseline responses (A/B) were randomized for each query.
    \item \textbf{Interaction:} Voters entered a clinical or theoretical query and received two simultaneous streams.
    \item \textbf{Metrics:} After reading both responses, voters could select "A is Better," "B is Better," or "Tie." Votes were stored in a JSONL format for statistical analysis.
\end{itemize}

\section{Quantitative Results}

The primary goal of the quantitative analysis was to determine if the TCM-Sage system provided a statistically significant improvement over a general-purpose LLM.

\subsection{Comparative Performance}
The distribution of the 56 votes is summarized below:
\begin{itemize}
    \item \textbf{TCM-Sage Wins:} 35 (62.5\%)
    \item \textbf{Plain LLM Wins:} 15 (26.8\%)
    \item \textbf{Ties:} 6 (10.7\%)
\end{itemize}

The TCM-Sage system outperformed the baseline in over 60\% of cases, more than doubling the win rate of the general-purpose search-augmented model.

\subsection{Statistical Significance}
To rigorously test the null hypothesis that there is no difference between the two systems, a \textbf{paired T-test} was performed on the win/loss data (assigning 1 for a win, -1 for a loss, and 0 for a tie).
\begin{itemize}
    \item \textbf{T-statistic:} 3.03
    \item \textbf{P-value:} 0.0037
\end{itemize}
With $p < 0.01$, the results are highly significant, allowing us to reject the null hypothesis.

\subsection{Effect Size}
\textbf{Cohen's d} was calculated to measure the magnitude of the improvement:
\begin{itemize}
    \item \textbf{Cohen's d:} 0.40
\end{itemize}
An effect size of 0.40 indicates a \textit{medium effect}, suggesting that while the general-purpose model is a capable baseline, the specialized TCM-Sage system provides a meaningful and consistent advantage in the quality of clinical reasoning and source accuracy.

\section{Temporal Analysis}

The evaluation was conducted over a four-day period (April 3 to April 6, 2026), during which the system underwent minor iterative refinements based on initial feedback.

\begin{itemize}
    \item \textbf{April 3:} 2 TCM-Sage / 1 Plain (67\% Win Rate) — Initial testing phase.
    \item \textbf{April 4:} 13 TCM-Sage / 8 Plain (62\% Win Rate) — Main testing batch from student groups.
    \item \textbf{April 5:} 8 TCM-Sage / 0 Plain (100\% Win Rate) — Testing following the integration of the \textit{Measurement Conversion} table.
    \item \textbf{April 6:} 12 TCM-Sage / 6 Plain (67\% Win Rate) — Extended testing with a wider group of practitioners.
\end{itemize}

The cumulative trend shows a consistent preference for TCM-Sage, with a notable peak in performance after specific clinical knowledge (e.g., unit conversions) was hard-coded into the system prompt.

\section{Qualitative Feedback and Case Studies}

Qualitative feedback from TCM professionals provided deeper insights into the specific strengths and weaknesses of both systems.

\subsection{Dosage and Unit Conversion}
A doctoral student at the Guangdong Provincial Hospital of TCM tested a query regarding the dosage of \textit{Mahuang Tang}. He observed that the general LLM (Baseline) correctly identified modern dosage ranges but the TCM-Sage system (at that time) provided a literal translation of Han Dynasty units that the model incorrectly converted (1 \textit{liang} $\approx$ 3g). 

\textit{Outcome:} This feedback led to the immediate implementation of the \textit{Measurement Conversion} table (1 \textit{liang} $\approx$ 13.8g) in the system prompt. Subsequent tests showed that practitioners found the corrected system far superior to the baseline, as it provided both historical accuracy and modern relevance.

\subsection{Source Traceability}
Students consistently highlighted the \textbf{Citation Panel} as the most valuable feature. One student noted: "The ability to see the original classical text immediately, without having to copy-paste the formula name into a search engine to verify the source, is the key value of this system." This indicates that the \textit{traceability} provided by RAG is as important as the \textit{content} of the generated response.

\subsection{Diagnostic Hedging and Safety}
A senior practitioner noted that the early version of the system was "over-assertive" in its diagnosis of \textit{Pi Xu} (Spleen Deficiency) based on limited symptoms. Following this, \textbf{Hedging Principles} were added to the prompt, requiring the system to use more tentative language ("The symptoms suggest a possibility of...") when data is incomplete. The practitioner later confirmed that the updated responses felt more like a professional consultation than a "black-box" diagnosis.

\subsection{Novel Clinical Insights}
In one recorded case, a student consulted the system for a sore throat. The system suggested \textit{Jiegeng Tang}, a classical formula the student had not considered. Upon reviewing the provided citations, the student confirmed the validity of the recommendation. This demonstrates the system's potential as a "clinical co-pilot" that can surface relevant but overlooked classical knowledge.

\section{Limitations of Evaluation}

While the results are statistically significant, several limitations must be acknowledged:
\begin{enumerate}
    \item \textbf{Sample Size:} 56 votes are sufficient for statistical significance but do not represent a comprehensive clinical trial.
    \item \textbf{Voter Demographic:} The majority of voters were TCM students rather than senior clinicians with decades of experience.
    \item \textbf{Evaluation Format:} The Arena tests overall user preference, which includes factors like formatting and tone, not just clinical accuracy.
    \item \textbf{Prompt Influence:} Because the system prompts differed between the two models, the evaluation measures the effectiveness of the \textit{entire pipeline} (Prompt + RAG + KG), not just the retrieval component in isolation.
\end{enumerate}

\section{Discussion of Results}

The results of the evaluation demonstrate that a specialized RAG system tailored for TCM significantly outperforms a general-purpose search-augmented LLM.

\begin{itemize}
    \item \textbf{Accuracy vs. Verifiability:} The primary advantage of TCM-Sage lies in its ability to provide verbatim classical citations. While the general LLM can often "summarize" TCM knowledge from the web, it frequently halluncinates formula origins or misinterprets classical syntax.
    \item \textbf{The Value of Domain Knowledge:} The integration of a specialized knowledge graph and a professional system prompt (including ancient-to-modern measurement conversions) is essential for clinical utility. Without these, even a high-quality RAG system is prone to fundamental errors in interpretation.
    \item \textbf{Impact on the Field:} The $p = 0.0037$ result strongly suggests that specialized AI tools can provide a higher standard of information for the TCM community than general-purpose assistants, particularly in the areas of citation traceability and classical methodology adherence.
\end{itemize}


\chapter{Discussion}

\section{Comparison with Existing TCM AI Systems}

TCM-Sage occupies a distinct position in the TCM AI landscape. While existing products pursue an ``AI Doctor'' paradigm, TCM-Sage functions as an evidence-synthesis reference tool. This distinction is not merely semantic; it reflects fundamentally different design philosophies with practical consequences.

General-purpose LLMs such as ChatGPT and Qwen operate as black boxes. When queried about TCM topics, they generate responses from parametric memory without any mechanism for source verification. A practitioner receiving a formula recommendation from ChatGPT cannot determine whether that recommendation originates from \textit{Shanghan Lun}, from a modern textbook, or from the model's interpolation of training data. In clinical contexts where incorrect dosage or formula selection carries real risk, this unverifiability is the primary obstacle to adoption, not hallucination per se.

Specialized TCM AI products have attempted to address domain accuracy but not transparency. HuatuoGPT achieved 88.1\% on TCM licensing examinations~\cite{huatuogpt2023}, and Qihuang 2.0 deployed a 320-billion-parameter multimodal model incorporating tongue and pulse image analysis. iFlytek's medical AI reached 107 primary care clinics with over 9,800 assisted diagnoses. Yet clinical adoption remains low: Qihuang reports less than 15\% usage among target practitioners. Published surveys identify three persistent barriers: data sovereignty concerns (hospitals refuse to upload patient data to cloud AI), liability ambiguity (who bears responsibility for AI-generated prescriptions), and the fundamental mismatch between AI-generated diagnoses and TCM's individualized \textit{Bian Zheng Lun Zhi} (Pattern Differentiation and Treatment) methodology.

TCM-Sage addresses these barriers through five design decisions:

\begin{enumerate}
    \item \textbf{Citation Traceability}: Every claim links to a specific classical text passage, visible through the Citation Panel. Practitioners can verify sources without copy-pasting text into external search tools.
    \item \textbf{Classical Text Specialization}: The 17-text corpus provides depth in foundational TCM theory that general LLMs cannot match, backed by clause-level retrieval for structured texts.
    \item \textbf{Full Local Deployment}: The entire pipeline, including embedding models and LLM inference, can run on consumer hardware via Ollama or LMStudio. No patient data leaves the institution's network.
    \item \textbf{Modular Knowledge Base}: Adding or removing knowledge sources requires only text ingestion, not model retraining. When practitioners suggested adding modern clinical texts during testing, the response was immediate: the architecture supports this without modification.
    \item \textbf{Ancient Measurement Accuracy}: An embedded conversion table provides historically verified unit conversions (Han dynasty \textit{liang} $\approx$ 13.8g, not the commonly misquoted 3g), preventing dosage calculation errors that other systems produce.
\end{enumerate}

The SymMap v2.0 knowledge graph deserves specific mention. Unlike the classical texts, SymMap is a modern pharmacological database integrating contemporary research. The crosswalk bridge connects these modern entities to classical terminology, creating a bidirectional link between historical medical knowledge and current pharmacological understanding. This hybrid of classical texts and modern knowledge graph is unique among TCM AI systems.

\section{Strengths Demonstrated}

The Arena evaluation provides quantitative evidence: with $p = 0.0037$ and Cohen's $d = 0.40$, the preference for TCM-Sage over the baseline is statistically significant at the 1\% level with a medium effect size. This result held consistently across four days of testing with different tester groups.

Qualitative feedback reinforced specific strengths. The Citation Panel received the strongest endorsement from practitioners, who described it as eliminating the need to manually verify AI-generated claims against source texts. One tester, a Year 5 TCM student, noted that the system recommended \textit{Jiegeng Tang} (Platycodon Decoction) for a sore throat case that he had not considered. Upon review, he confirmed this as a valid and potentially superior approach from a classical perspective. This case demonstrates that RAG-based evidence synthesis can surface treatment options that human practitioners might overlook due to training focus or cognitive bias.

The modular architecture received indirect validation when practitioners at Guangdong Provincial Hospital of TCM suggested expanding the corpus to include modern clinical masters' works (Zhang Xichun's \textit{Yi Xue Zhong Zhong Can Xi Lu}, Shi Xuemin's acupuncture texts, and Dong's Extraordinary Points system). The fact that such expansion is architecturally straightforward confirms that the design decision to use RAG over fine-tuning was sound.

\section{Limitations and Honest Assessment}

Several limitations emerged during testing and must be acknowledged.

\textbf{Knowledge Boundary Errors.} The system's accuracy is bounded by its corpus. When a practitioner asked about \textit{Mahuang Tang} dosage, the system stated that one Han dynasty \textit{liang} equals approximately 3 grams, a significant error (the correct value is 13.8--15.6g based on archaeological evidence). This occurred because the classical texts contain dosage ratios but not modern unit conversions. The plain LLM baseline, with access to web search, retrieved the correct conversion. This limitation was addressed by embedding a historically verified measurement conversion table directly into the system prompt.

\textbf{Corpus Coverage Gaps.} The 17-text corpus covers foundational theory, \textit{Shanghan}/\textit{Wenbing}, herbology, formulary, and acupuncture. It does not cover surgery, orthopedics, ophthalmology, specialized gynecology or pediatrics texts, or pulse diagnosis monographs. Practitioner feedback was direct: ``These texts are too traditional. Real clinical learning requires modern clinical masters' experience and knowledge systems.'' This feedback, while critical, validates the system's modular design as the appropriate architecture for iterative corpus expansion.

\textbf{Diagnostic Over-Assertion.} The system initially presented diagnostic interpretations with excessive confidence. When analyzing a pediatric case, it concluded \textit{Pi Xu} (spleen deficiency) without sufficient evidence in the case data, and classified an initial rash presentation as \textit{Re Yu Yu Ying} (heat trapped in the nutritive level) when the clinical presentation better supported \textit{Xie Qi Zai Biao} (pathogen at the exterior). Hedging principles were added to the system prompt to mitigate this tendency, but the underlying challenge of calibrating LLM confidence in medical reasoning remains partially unresolved.

\textbf{Evaluation Scale.} The 56-vote Arena sample, while statistically sufficient for significance testing, represents a limited evaluator pool. Most voters were TCM students rather than senior practitioners. The Arena format measures overall preference but does not decompose quality into specific dimensions such as citation accuracy, diagnostic reasoning, or readability.

\section{Ethical Considerations}

TCM-Sage is positioned as an evidence reference tool, not a diagnostic system. This distinction carries ethical weight. The system does not claim to diagnose, and its outputs carry an explicit disclaimer that all suggestions are for clinical reference only, with final diagnostic and prescriptive authority belonging to licensed practitioners.

Data privacy is addressed architecturally. Full local deployment eliminates the need to transmit patient data to external servers. During testing, practitioners at Guangdong Provincial Hospital of TCM noted this as a meaningful advantage, given institutional policies restricting patient data sharing with cloud AI services.

The use of AI tools in developing this project is disclosed in the Acknowledgments section, in compliance with university policy.

\begin{figure}[H]
\centering
\fbox{\parbox{0.8\textwidth}{\centering [PLACEHOLDER: Comparison table --- TCM-Sage vs existing TCM AI systems across key dimensions (citation traceability, deployment model, knowledge base type, clinical approach)]}}
\caption{Feature comparison between TCM-Sage and existing TCM AI systems.}
\label{fig:comparison}
\end{figure}


\input{chapters/chapter7}


\chapter*{Acknowledgments}
\addcontentsline{toc}{chapter}{Acknowledgments}

I would like to express my sincere gratitude to my supervisor, Dr.\ Zhang Ce, for his guidance throughout this project. His suggestions on evaluation methodology---specifically the use of paired t-tests and Cohen's $d$ for statistical validation---and his discussions on blind A/B experimental design were instrumental in shaping the Arena evaluation framework.

I am deeply grateful to Kenny Woo Shi Nam, a Year 5 student at the HKBU School of Chinese Medicine (cGPA 3.97/4.0), currently interning at Guangdong Provincial Hospital of TCM. Kenny's contributions were extensive: he advised on the selection of the 17 classical texts, conducted rigorous testing across formula retrieval, clause comparison, and clinical case analysis, and provided detailed feedback on diagnostic accuracy and system usability that directly shaped multiple system improvements.

I also thank the doctoral students at Guangdong Provincial Hospital of TCM who participated in the Arena evaluation and provided clinical feedback, including specific recommendations for corpus expansion and insights on bridging classical theory with modern clinical practice.

AI-powered development environments (OpenCode, Cursor) were used for code implementation assistance throughout the project, with all code reviewed and tested by the author. AI writing tools were used for grammar and style refinement during report preparation. All architectural decisions, system design, and written content are the author's own work.

\appendix
\chapter{Classical Text Corpus}

Table~\ref{tab:corpus_detail} provides the complete list of classical TCM texts included in TCM-Sage, with their dynasty of origin, attributed author, chunking strategy, and resulting chunk count.

\begin{table}[H]
\centering
\caption{Detailed corpus breakdown of the 17 classical texts.}
\label{tab:corpus_detail}
\begin{tabularx}{\textwidth}{lllrr}
\toprule
\textbf{Text} & \textbf{Dynasty} & \textbf{Author} & \textbf{Chunks} & \textbf{Strategy} \\
\midrule
Bencao Gangmu & Ming & Li Shizhen & 5,169 & Chapter \\
Beiji Qianjin Yaofang & Tang & Sun Simiao & 2,326 & Chapter \\
Rumen Shiqin & Jin & Zhang Congzheng & 920 & Chapter \\
Jingui Yaolue & E. Han & Zhang Zhongjing & 489 & Clause \\
Danxi Xinfa & Yuan & Zhu Zhenheng & 413 & Chapter \\
Shennong Bencao Jing & W. Han & Unknown & 401 & Chapter \\
Shanghan Lun & E. Han & Zhang Zhongjing & 388 & Clause \\
Zhenjiu Jiayijing & W. Jin & Huangfu Mi & 376 & Chapter \\
Lanshi Micang & Yuan & Li Dongyuan & 318 & Chapter \\
Wenbing Tiaobian & Qing & Wu Jutong & 313 & Chapter \\
Huangdi Neijing Suwen & W. Han & Unknown & 295 & Chapter \\
Huangdi Neijing Lingshu & W. Han & Unknown & 227 & Chapter \\
Suwen Xuanji Yuanbing Shi & Jin & Liu Wansu & 205 & Chapter \\
Piwei Lun & Yuan & Li Dongyuan & 158 & Chapter \\
Neiwaishang Bianhuo Lun & Yuan & Li Dongyuan & 95 & Chapter \\
Bashiyi Nanjing & E. Han & Unknown & 88 & Chapter \\
Wenre Lun & Qing & Ye Tianshi & 23 & Chapter \\
\midrule
\textbf{Total} & & & \textbf{12,204} & \\
\bottomrule
\end{tabularx}
\end{table}

\chapter{System Setup and Reproduction Guide}

This appendix describes how to set up and run TCM-Sage from the source repository.

\section*{Prerequisites}
\begin{itemize}
    \item Python 3.11+ with \texttt{pip}
    \item Node.js 18+ with \texttt{npm}
    \item A DashScope API key (or alternative LLM provider credentials)
    \item Git
\end{itemize}

\section*{Step 1: Clone and Install}
\begin{verbatim}
git clone https://github.com/AndyZHENG0715/TCM-Sage.git
cd TCM-Sage

# Python backend
python -m venv venv
venv\Scripts\activate          # Windows
pip install -r requirements.txt

# Frontend
cd web && npm install && cd ..
\end{verbatim}

\section*{Step 2: Configure Environment}
Copy \texttt{.env.example} to \texttt{.env} and fill in the required API keys:
\begin{verbatim}
DASHSCOPE_API_KEY=your_key_here
LLM_PROVIDER=alibaba
\end{verbatim}
See \texttt{docs/CONFIG.md} for the full list of configurable environment variables.

\section*{Step 3: Build the Vector Index}
\begin{verbatim}
venv\Scripts\python src/ingest.py
\end{verbatim}
This processes the 17 classical texts in \texttt{data/source/}, generates embeddings via DashScope, and stores them in ChromaDB at \texttt{vectorstore/chroma/}. The process supports checkpoint and resume. Initial ingestion takes approximately 30--60 minutes due to API batch limits.

\section*{Step 4: Start the Application}
\begin{verbatim}
# Terminal 1: Backend (port 8000)
venv\Scripts\python src/api.py

# Terminal 2: Frontend (port 3000)
cd web && npm run dev
\end{verbatim}
Access the application at \texttt{http://localhost:3000}.

\section*{Step 5: Local Deployment (Optional)}
For fully local operation without cloud API calls, install Ollama and configure:
\begin{verbatim}
LLM_PROVIDER=ollama
OLLAMA_BASE_URL=http://localhost:11434/v1
\end{verbatim}
Note: the current deployment uses DashScope cloud embeddings. For fully offline operation, a local embedding model must be substituted in \texttt{src/embeddings.py}.


\bibliographystyle{abbrvnat}
\bibliography{references}

\end{document}
